\chapter{Research Methodology}
\label{chap:analysis}

This chapter presents the research methodology applied in this thesis. It describes 
the overall research design, the operationalization of the Integration Maturity Model, 
the design of the survey instrument and expert interviews, the data collection 
procedure, and the statistical methods used to analyze the collected data.

\section{Research Design and Approach}
\label{sec:design}

This thesis adopts a quantitative research design, supplemented by qualitative 
expert interviews. The primary research strategy is a cross-sectional survey study, 
in which data is collected from multiple organizations at a single point in time. 
This design is appropriate for the research aim, which is to assess the current 
distribution of IS integration maturity levels across PMO environments in Latvian 
enterprises and examine its relationships with reporting automation, data quality, 
and governance standardization \cite{saunders2019}.

The philosophical stance of the study is post-positivist, assuming that an 
objective reality exists and can be measured through structured instruments, 
while acknowledging that measurement is subject to some degree of error. This 
stance is consistent with the use of Likert-scale survey items and statistical 
analysis as the primary data collection and analysis methods.

The research follows a deductive approach, in which the four hypotheses stated 
in the Introduction are derived from the theoretical framework presented in 
Chapter~1 and tested against empirical data collected from Latvian enterprises.

\section{Operationalization of the Integration Maturity Model}
\label{sec:operationalization}

The four-level Integration Maturity Model (IMM) proposed in Chapter~1 serves as 
the primary measurement framework. To enable survey-based measurement, each 
maturity level is operationalized through a set of observable, measurable 
indicators, as presented in Table~\ref{tab:imm-operationalization}.

{\singlespacing
\begin{table}[h]
  \caption{Operationalization of Integration Maturity Model levels}
  \label{tab:imm-operationalization}
  \centering
  \begin{tabular}{p{0.08\linewidth}p{0.15\linewidth}p{0.35\linewidth}p{0.30\linewidth}}
    \toprule
    \textbf{Level} & \textbf{Name} & \textbf{Observable Indicators} & 
    \textbf{Survey Items}\\
    \midrule
    1 & Initial & No system connections; all data consolidated manually in 
    spreadsheets; no automated reports & 
    "Data is consolidated manually from multiple tools into spreadsheets"\\
    2 & Developing & Some file-based exports between systems; partial 
    automation; scheduled report generation & 
    "Some systems share data via file exports (CSV/Excel)"\\
    3 & Defined & API connections between key systems; automated data flows; 
    near real-time reporting & 
    "Systems are connected via APIs; reports are generated automatically"\\
    4 & Optimized & Full platform integration; unified BI dashboard; 
    real-time portfolio visibility & 
    "All systems are fully integrated in a unified platform with real-time dashboards"\\
    \bottomrule
  \end{tabular}
\end{table}
}

Respondents are asked to select the description that best matches their 
organization's current PMO IS integration approach. The selected level serves 
as the primary independent variable in the analysis.

\section{Survey Instrument Design}
\label{sec:survey}

The survey instrument was designed to measure the key constructs of the study: 
integration maturity level, reporting automation, data quality, and governance 
standardization. The instrument consists of four parts:

\begin{itemize}
    \item \textbf{Part A --- Organizational Profile:} Questions about the 
    organization's size, industry, and PMO structure (categorical items).
    \item \textbf{Part B --- Integration Maturity Assessment:} A single 
    categorical item in which respondents select the IMM level that best 
    describes their PMO's IS integration approach, supplemented by 
    Likert-scale items measuring specific integration practices.
    \item \textbf{Part C --- Reporting Automation and Data Quality:} 
    Five-point Likert-scale items measuring the degree of reporting 
    automation and perceived data accuracy and timeliness.
    \item \textbf{Part D --- Governance Standardization:} Five-point 
    Likert-scale items measuring the degree of governance standardization 
    in PMO processes.
\end{itemize}

All Likert-scale items use a five-point response scale ranging from 
1 (Strongly Disagree) to 5 (Strongly Agree). The survey was designed in 
English and distributed electronically via Google Forms. The estimated 
completion time is 10--15 minutes.

\section{Expert Interview Design}
\label{sec:interviews}

To supplement the quantitative survey findings with qualitative depth, 
2--3 structured expert interviews are conducted with PMO leads, IT architects, 
or digital transformation managers from Latvian enterprises. The interviews 
follow a structured protocol with open-ended questions covering the following 
themes:

\begin{itemize}
    \item The current state of IS integration in the respondent's PMO environment
    \item Challenges experienced due to low integration maturity
    \item The role of BI and dashboard systems in governance standardization
    \item Perceived barriers to advancing integration maturity
    \item Practical recommendations for improving PMO IS integration
\end{itemize}

Interview responses are analyzed qualitatively using thematic analysis to 
identify patterns and insights that complement the quantitative findings.

\section{Data Collection Procedure}
\label{sec:datacollection}

The survey is distributed to representatives of medium and large Latvian 
enterprises (50+ employees) with formal PMO structures. The target sample 
size is 30--50 responses, which is sufficient for the planned descriptive 
and correlational statistical analyses \cite{field2018}. Respondents are 
identified through professional networks, LinkedIn, and direct organizational 
contacts. Only one response per organization is accepted to ensure independence 
of observations.

Data collection takes place over a period of four to six weeks. All responses 
are anonymous and stored securely in accordance with data protection principles. 
Participation is voluntary and respondents are informed of the purpose of the 
study prior to completing the survey.

\section{Data Analysis Methods}
\label{sec:analysis}

The collected data is analyzed using the following statistical methods:

\textbf{Descriptive statistics.} Frequency distributions and percentages are 
used to describe the distribution of integration maturity levels across the 
sample. Mean scores and standard deviations are calculated for Likert-scale 
items measuring reporting automation, data quality, and governance 
standardization.

\textbf{Normality testing.} The Shapiro-Wilk test is applied to determine 
whether the data follows a normal distribution, which informs the choice 
between parametric (Pearson) and non-parametric (Spearman) correlation 
analysis \cite{field2018}.

\textbf{Correlation analysis.} Pearson or Spearman correlation coefficients 
are calculated to examine the relationships between integration maturity level 
and the three dependent variables: reporting automation, data quality, and 
governance standardization. Statistical significance is assessed at the 
$p < 0.05$ level.

\textbf{Regression analysis.} Simple linear regression is applied to examine 
the predictive relationship between integration maturity level and each 
dependent variable, allowing for quantification of the strength and direction 
of the relationships.

\textbf{Qualitative analysis.} Expert interview transcripts are analyzed using 
thematic analysis to identify recurring themes and insights that supplement 
the quantitative findings.

All quantitative analyses are conducted using SPSS or R statistical software.

\section{Chapter Summary}
\label{sec:chap2summary}

This chapter described the research methodology applied in this thesis. 
A quantitative cross-sectional survey design with qualitative expert interview 
supplementation was adopted. The four-level Integration Maturity Model was 
operationalized into measurable survey indicators. The survey instrument was 
designed to measure integration maturity, reporting automation, data quality, 
and governance standardization across PMO environments in Latvian enterprises. 
Data collection targets 30--50 responses from medium and large Latvian 
enterprises with formal PMO structures. Statistical analysis methods including 
descriptive statistics, normality testing, correlation analysis, and regression 
analysis are applied to test the four research hypotheses. The results of the 
empirical study are presented in Chapter~3.
